\section{Problem tanımı}\label{sec:problem}

$P=\{v_1,\dots,v_n\}\subset\mathbb{R}^2$ hedefleri (ağaçlar, delme noktaları,
muayene noktaları), $D\in\mathbb{R}^2$ bir depoyu (baz istasyonu), $k$ ise özdeş robot
sayısını göstersin. Bir çözüm, $P$'nin $k$ adet (boş olabilen) $R_1,\dots,R_k$ sıralı
dizisine bölüntüsüdür; $j$ numaralı robot $D$'den ayrılır, $R_j$ dizisindeki hedefleri sırayla ziyaret eder ve
$D$'ye geri döner. Öklidyen uzaklık $d(\cdot,\cdot)$ olmak üzere, bir
$R=(u_1,\dots,u_m)$ rotasının uzunluğu
\begin{equation}
L(R) = d(D,u_1) + \sum_{t=1}^{m-1} d(u_t,u_{t+1}) + d(u_m,D),
\end{equation}
ve amaç \emph{tamamlanma süresi}
\begin{equation}
\min_{R_1,\dots,R_k}\; C_{\max} = \max_{j} L(R_j),
\end{equation}
olup bu problem, tek depolu min--maks çoklu gezgin satıcı problemidir
(min--maks mTSP) \citep{bektas2006,necula2015}. Robot hızı sabit $v$ olduğunda
filo turu $C_{\max}/v$ sürede tamamlanır. Filonun tümü her turda birlikte planlanıp
birlikte yola çıkıyorsa ve planlama ile yürütme birbirini izliyorsa, $r$ eşzamanlı
turun toplam süresi
\begin{equation}\label{eq:mission}
T_{\text{mission}} = \sum_{\rho=1}^{r}\Big(t_{\text{plan}}^{(\rho)} + C_{\max}^{(\rho)}/v\Big)
\end{equation}
olur. Bu denklem eşzamanlı tur modelinin özel durumudur; aşağıdaki olay tabanlı
hasat benzetiminde görev süresi, planlama ve çalışma zamanları örtüşebildiği için
bu toplamla değil, son platformun bitiş olayının zamanıyla tanımlanır.
Tüm uzaklıklar gerçel değerli alınmıştır (TSPLIB yuvarlaması uygulanmamıştır); LKH-3 sonuçları da dâhil olmak üzere
raporlanan her tamamlanma süresi, döndürülen rotalardan çift duyarlıkla yeniden hesaplanmıştır.

\paragraph{Alt sınır} Her rota $D$'den geçen kapalı bir yürüyüş olduğundan, uzunluğu
en az köşelerinin en küçük kapsayan ağacının (MST) uzunluğu kadardır,
$L(R_j)\ge \mathrm{MST}(R_j\cup\{D\})$; bu $k$ ağacın birleşimi ise $P\cup\{D\}$'nin
bağlantılı bir kapsayan alt çizgesidir. Ayrıca her robot en uzak hedefine
ulaşıp geri dönmek zorundadır. Dolayısıyla
\begin{equation}\label{eq:lb}
C_{\max}^{*} \;\ge\; \mathrm{LB} = \max\Big\{\,2\max_{v\in P} d(D,v),\;\; \mathrm{MST}(P\cup\{D\})/k\,\Big\},
\end{equation}
olup bu sınır Delaunay üçgenlemesi yardımıyla $O(n\log n)$ sürede hesaplanır. Sınır yalnızca
mutlak kalite için bir tutarlılık denetimi olarak kullanılmıştır; küçük $k$ değerleri için sıkı değildir.

\paragraph{Hasat varyantı} Hasat uygulamasında
hedefler, bir hizmet süresi (olgun meyvenin toplanması) gerektiren ağaçlardır; her robot,
kapasitesi $Q$ olan tek bir kasa taşıyan ve kasa istasyonuna birkaç kez dönebilen
bir hasat platformudur; hareket sıra aralarıyla sınırlıdır ve ağaç başına olgun
meyve sayısı ağaç gözlemlenene kadar yalnızca yaklaşık olarak bilinir.
Amaç, son platformun işini bitirdiği $C_j$ zamanlarının en büyüğü,
$T_{\text{mission}}=\max_j C_j$ olmaktadır. Bir platform $t$ zamanında döndüğünde,
planlamaya $s\ge t$ zamanında başlanabiliyorsa yeni kalkış zamanı
$\max\{t+t_{\text{bin}},\,s+t_{\text{plan}}\}$ olur; burada $t_{\text{bin}}$
kasa değişim süresidir. Böylece planlama diğer platformların çalışmasıyla ve
kasa değişimiyle örtüşür; görev süresi yalnızca kritik kalkış gecikmelerini içerir.
Bu amaç olay tabanlı benzetimle değerlendirilmiştir.
